%% kapa: exam=1 round=19 year=2008 period=2 form=A subject=english subject_ko=영어 seq=21 number=61
%% source: Q-Net 감정평가사 기출문제 (cst00309) artlSeq=1001962 file=extracted/1008/1001962/2008년 제19회 감정평가사 1차시험 기출문제/2교시.hwp
%% conversion: kordoc 4.24.51 -> tools/gen1.py (auto); review=proofread-sonnet; flags=none
\begin{question}{61}
% [공통 지문/지시문: 원문에서 여러 문항에 걸쳐 제시됨]
\begin{framedtext}
※ 빈칸에 들어갈 가장 알맞은 표현을 고르시오. (60-62) \par
\end{framedtext}
Now it seems that \blank{} keeps marriages between equals is not rigid formulas of fairness, but tolerance for change, flexibility, and negotiation.
\begin{choices}
  \item how
  \item what
  \item it
  \item whichever
  \item which
\end{choices}
\end{question}
\answer{2}
